\begin{table}[H]
\caption{Distributional diagnostics of the Diebold--Mariano comparisons. $n$ is the number of aligned test-set forecast errors per unit; IQR is the interquartile range of the DM statistic; loss differences are mean standard loss minus mean fractional loss (positive values favour the fractional model); FDS is the fractional dominance score (significant fractional wins divided by the number of units with a defined test); ``Majority'' names the model with the larger number of significant wins. The mean loss difference is dominated by a few high-incidence units and is reported for completeness only.\label{tab:dm-diagnostics}}
\begin{adjustwidth}{-\extralength}{0cm}
\small
\begin{tabularx}{\fulllength}{lCCCCCCCC}
\toprule
\textbf{Comparison} & \textbf{Median $n$} & \textbf{Median DM} & \textbf{IQR of DM} & \textbf{Median $p$} & \textbf{Median loss diff.} & \textbf{Mean loss diff.} & \textbf{FDS} & \textbf{Majority}\\
\midrule
ARIMA vs. ARFIMA & 229 & -5.855 & -23.086 to 6.684 & \mbox{$<10^{-16}$} & -288.82 & \mbox{$1.10 \times 10^{10}$} & 0.333 & standard\\
RF vs. Fractional RF & 221 & -1.288 & -2.786 to 1.542 & 0.029 & -137.88 & \mbox{$-2.52 \times 10^{6}$} & 0.203 & standard\\
XGBoost vs. Fractional XGBoost & 221 & 0.258 & -2.022 to 2.466 & 0.020 & 8.20 & \mbox{$1.65 \times 10^{7}$} & 0.276 & fractional\\
\bottomrule
\end{tabularx}
\end{adjustwidth}
\end{table}
